% Pursuing Stacks (cotes 134-2 to 134-8): archive cote 134-6, site page 85 (PDF page 86), typescript page 403.
% typescript. Draft compared: ps5.tex:5801-5864.
% AI-transcribed by Claude Opus 5.5 (claude-opus-5-5, adaptive thinking, effort high), Message Batches API, batch ps3 (msgbatch_01XdFmuVVBAcgc7EFxdtCqsj), 2026-11-12; not yet checked by eye.
\scanpage{134-6/85}
\mspage{403}

Associating to any small category P the algebraic abelian topos
\uline{Homadd}(\unsure{$P^{o}$;}~(Ab)), should be viewed of course as the abelian \unsure{analogon|analggon} of
$A \hand{$\mapsto$} A^{\wedge}$, associating to a small category A the corresponding algebraic
topos . It merits a notation of it's own, s a y
\[ P^{\&} \overset{\mathrm{def}}{=} \uline{Homadd}\ (P^{o}, (Ab)) \ , \]
and as \xout{the} in the non additive case, we get a canonical inclusion functor
\[ \hand{$\varepsilon$}_{P} \ : \ P \longrightarrow P^{\&} \]
which is additive\hand{.} \hand{(}\over{i}{I}t's composition with the canonical functor $P^{\&}\hand{$\to$} P^{\hand{$\wedge$}}$ is
\xout{Moreover} \xout{being}the canonical inclusion previously denoted by $\hand{$\varepsilon$}_{P}$ too from P to $P^{\wedge}$.\hand{)}\hand{.}
Next thing we'll expect, in analogy to prop.3 , is that \xout{the} for any
additive category \hand{$\mathcal{M}$} stable under small direct limits, the following
canonical functor is an equivalence of categories:
\margin{(**)}
\[ \uline{Hom}_{\hand{$!$}}(P^{\&}, \hand{$\mathcal{M}$}) \hand{$\xrightarrow{\approx}$} \uline{Homadd}(P, \hand{$\mathcal{M}$}) \ , \quad F \longmapsto F \circ \hand{$\varepsilon$}_{P} \ . \]
T\inserted{h}e proof, via construction of a quasi-inverse functor, should be about
the same as for prop.3, which should go through once we get the abelian
analagon of the well-known fact in $A^{\wedge}$ , that any object\inserted{F} in $A^{\wedge}$ can be
recovered as direct limit in $A^{\wedge}$\xout{\illegible{1}} of objects of $A$, according to $A_{/F}$ as
an indexing category - which makes us expect that we get \hand{too} :
\[ F \hand{$\xleftarrow{\sim}$} \underset{a \text{ in } P^{\&}_{/F}}{\lim_{\hand{$\to$}}} a \qquad (\text{direct limit in } P^{\&}) \ . \]
\inserted{From (**),}
We \deleted{then} get as in cor.1, passing to the dual categories, the dual equi-%
valence
\[ \uline{Hom}^{\hand{$!$}}(P^{\& o} \ , \hand{$\mathcal{N}$}) \hand{$\xrightarrow{\approx}$} \uline{Homadd}(P^{o}, \hand{$\mathcal{N}$}) \ , \]
valid if \hand{$\mathcal{N}$} is an additive category stable under small inverse limits,
\deleted{and} \over{h}{H}ence, if \hand{$\mathcal{M}$} is \inserted[typed]{additive and} stable under both types of limits, the equivalen\unsure{ce}
\margin{(***)}
\[ \uline{Hom}_{\hand{$!$}}(P^{\&}, \hand{$\mathcal{M}$}) \hand{$\xrightarrow{\approx}$} \uline{Hom}^{\hand{$!$}}(Q^{\& o}, \hand{$\mathcal{M}$}) \ , \quad \text{with } Q{=}P^{o} \ , \]
between abelian cosheaves on $P^{\&}$ and abelian sheaves in $Q^{\&}$ , with values
in the same \xout{category} additive category (in analogy to (39)). In the
particular case \hand{$\mathcal{M}$}= (Ab), this then gives rise to the equivalence
\[ Q^{\&} \hand{$\simeq$} \uline{Hom}_{\hand{$!$}}(P^{\&}, (Ab)) \]
and to the corresponding pairing
\[ P^{\&} \ x \ Q^{\&} \longrightarrow (Ab) \]
which is a duality, namely induces an equivalence between each \deleted{category}\inserted{abelian topos}

\margin{85}
